\title[Single-centred indices of the $\ZZ_2$ CHL model]{Every Single-Centred Black Hole Index\\ of the $\mathbb Z_2$ CHL Model Is Positive}

\author{Deep Bhattacharjee}
\address{Formerly, Electro-Gravitational Space Propulsion Laboratory (EGSPL), Bhubaneswar, Odisha 751030, India}
\address{Formerly, University of Calcutta (Heramba Chandra College), 23/49 Gariahat Road, Golpark, Kolkata 700029, West Bengal, India}
\email{d.bhattacharjee@erl-forschung.de}
\email{itsdeep@live.com}
\email{egspl@proton.me (GitHub-related contact; GitHub account creelie)}
\urladdr{https://orcid.org/0000-0003-0466-750X}

\begin{abstract}
Every single-centred black hole index of the $\ZZ_2$ CHL model is positive. The proof rests on an exact Rademacher expansion at level two.

\medskip\noindent\textbf{Keywords:} CHL models; single-centred black holes; mock Jacobi forms.

\smallskip\noindent\textbf{MSC 2020:} 11F50 (primary); 11F37, 11F30, 81T30, 83E30.

\smallskip\noindent\textbf{PACS:} 11.25.-w; 04.70.Dy; 11.25.Mj; 02.10.De.
\end{abstract}

\maketitle

\section{Introduction}\label{sec:intro}

A quarter-BPS dyon of $\mathcal N=4$ string theory in four dimensions carries an electric charge $Q$ and a magnetic charge $P$, and the protected index that counts such states is the helicity trace $B_6$. When the charges admit a single-centred black hole whose only hair are the zero modes of the broken supersymmetries, the twelve fermion zero modes fix the sign of the trace, and the horizon degrees of freedom contribute $-B_6>0$. Sen turned this observation into a conjecture about Fourier coefficients of Siegel modular forms \cite{Sen11}: the index $d^*(Q,P)$ of the single-centred black hole, extracted from the inverse of the dyon partition function with the contour of the attractor region, is positive.\footnote{Sen checked the sign on the charges of his Tables~1 to~5, for the heterotic string on $T^6$ and for the CHL models of orders $2,3,5,7$, and stated the conjecture for all charges in the cones (3.7)--(3.12) of \cite{Sen11}. The entropy function and attractor background behind the statement are reviewed in \cite{Sen08}.}

For type~II string theory on $K3\times T^2$, where the partition function is the Igusa cusp form $\Phi_{10}$ \cite{DVV97,SSY06,GN97}, the conjecture is settled. Bringmann and Murthy proved it for the two lowest magnetic charges $P^2=2,4$ by writing the finite parts of the first two Fourier--Jacobi coefficients of $1/\Phi_{10}$ through class numbers \cite{BM13}. Rossell\'o then treated all charges, by writing the single-centred index as a mixed Rademacher series and showing that one of its terms dominates the others; the existence of such series on $SL_2(\ZZ)$ is due to Bringmann and Manschot \cite{BMa13}, and Ferrari and Reys used them for these black holes \cite{FR17}. The CHL models \cite{CHL95} are the $\ZZ_N$ orbifolds of type~IIB on $K3\times T^2$ in which the $\ZZ_N$ acts on $K3$ by a symplectic automorphism and on $T^2$ by a shift; their partition functions $\Phit_k$ are Siegel modular forms of weight $k$ for subgroups of $Sp(2,\ZZ)$ \cite{JS06,DJS06,DS06,DN07}.\footnote{For $N=2$ the model is the heterotic string on $T^6$ obtained from the $E_8\times E_8$ string by exchanging the two $E_8$ factors together with a half shift on a circle \cite{CHL95}; for the ten-dimensional $E_8\times E_8$ theory and its relation to M-theory see \cite{SBBT26}. Background on the involutions of $K3$ surfaces and on Calabi--Yau threefolds fibred by $K3$ surfaces, which underlie the dual type~II description, is in \cite{BN26,Bha22}.} Bhand, Sen and Singh showed that the Fourier--Jacobi coefficients of $1/\Phit_k$ are meromorphic Jacobi forms, split them into an Appell--Lerch polar part and a mock Jacobi finite part, and identified the single-centred index with the Fourier coefficients of the finite part on an explicit set of charges \cite[Theorem~5.4]{BSS25}. Cardoso, Nampuri and Rossell\'o computed the polar coefficients of these finite parts through continued fractions and wrote the immortal index as a Rademacher series, taking the existence of such an expansion at level $N$ as an assumption \cite{CNR21}.

This paper settles the $\ZZ_2$ CHL model. It has $k=6$, duality group $\Gamma_1(2)$, and charges normalised so that $m=Q^2$, $n=P^2/2$ and $\ell=Q\cdot P$ are integers; the discriminant is $D=Q^2P^2-(Q\cdot P)^2=2mn-\ell^2$.

\begin{mainthm}\label{thm:all}
In the $\ZZ_2$ CHL model the single-centred index $d^*(Q,P)$ is positive for every charge with $D>0$.
\end{mainthm}

The proof has two halves. The first is an exact formula for the coefficients of the finite parts.

\begin{mainthm}\label{thm:radB}
For every $m\ge1$ the coefficients $c^F_m(n,\ell)$ of the finite part of the $m$-th Fourier--Jacobi coefficient of $1/\Phit_6$ with $D=2mn-\ell^2>0$ are given by an absolutely convergent vector-valued Rademacher series over the three cosets of $\Gamma^0(2)$ in $SL_2(\ZZ)$, with Bessel terms $I_{15/2}$ from the principal parts and correction terms $I_8$ and $I_{17/2}$ from the shadow. Its principal parts at the cusps $0$ and $\infty$ are explicit sums over Farey walls.
\end{mainthm}

The precise statement is Theorem~\ref{thm:rad}. The wall sums at the cusp $0$ have the shape of the continued-fraction formula of Cardoso, Nampuri and Rossell\'o; we prove them, and their counterpart at the cusp $\infty$, by following the Fourier coefficients of $1/\Phit_6$ across the walls of its poles, which are the Farey geodesics with an even denominator (Figure~\ref{fig:chambers}), and of a second product at the cusp $0$ across all Farey geodesics. The expansion itself comes from the circle method of Rademacher and Zuckerman \cite{RZ38,Zuc39}, applied to the theta coefficients of the finite parts once the Appell--Lerch sum has been completed in the manner of Zwegers \cite{Zwe02}; the completion is proved directly, by a Mordell integral, and the transformation law in all three frames follows from the product formula alone, without any statement about Siegel modular forms.

The second half is the sign. Write $x=\sqrt D/m$ and $\theta=\pi\ell/m$, and let $\mathcal E(D,m)$ be the contribution of the most polar term at the cusp $0$, an explicit multiple of a Bessel function $I_{15/2}$. Every charge has, in its duality orbit or after the Fricke exchange of $Q^2$ and $P^2/2$, a representative with $m\le n$ and $0\le\ell\le m$ (Proposition~\ref{prop:reduction}); then $x\ge1$. The terms of the Rademacher series attached to the leading family of Farey walls, those with the matrix $\begin{psmallmatrix}1&0\\1&1\end{psmallmatrix}$, add up, after division by $\mathcal E$, to an approximation of
\[
|P(z)|^2,\qquad P(z)=\prod_{n\ge1}(1-z^n)^{-8}(1-z^{2n})^{-8},\qquad z=e^{-\pi x+i\theta},
\]
the squared modulus of the generating function of the degeneracies of the electric half-BPS states of the model (Theorem~\ref{thm:model}). The function $|P|^2$ stays above $0.538$ on the disc $|z|\le e^{-\pi}$, its minimum being attained at $z=-e^{-\pi}$ (Proposition~\ref{prop:pmin}, Figure~\ref{fig:pdisc}). Everything else, namely the error of the approximation, the remaining walls, the other two cosets, the terms with $\gamma\ge2$ and the shadow corrections, is bounded by $0.347$ for $Q^2\le100$ by a certificate in interval arithmetic and by $0.3382$ for $Q^2\ge101$ by explicit inequalities (Theorems~\ref{thm:cert} and~\ref{thm:analytic}). The remaining charges have $D\le\Dexact$, and for them the index is computed exactly from the product formula (Theorem~\ref{thm:bounded}).\footnote{For $SL_2(\ZZ)$ one dominant term suffices. At level two the dominant contribution is a whole family, and the family resums to an eta quotient; the bound $|P|^2\ge0.538$ is sharp in the limit (Remark~\ref{rem:sharp}).}

On the first row the theorem also has a short proof by closed formulas. Following Bringmann and Murthy, an identity between Jacobi forms puts $\varphi_1$ into closed form through Hurwitz class numbers and the theta series of $D_4$, and the sign of every coefficient can be read off (Theorem~\ref{thm:Q1}).

The proof draws throughout on the work of Sen on the positivity conjecture and the attractor mechanism, of Bhand, Sen and Singh on the mock modularity of the CHL partition functions, of Cardoso, Nampuri and Rossell\'o on the decay walls, of Bringmann and Manschot and of Ferrari and Reys on mixed Rademacher series, on the circle method of Rademacher and Zuckerman, on the completions of Zwegers and on the theory of Jacobi forms of Eichler and Zagier \cite{EZ85}.

Section~\ref{sec:setup} recalls the partition function and the decomposition of Bhand, Sen and Singh, and Section~\ref{sec:Q1} treats the first row. Section~\ref{sec:reduction} proves the reduction and the bounded-discriminant theorem, and Section~\ref{sec:walls} describes the Farey walls. Section~\ref{sec:rademacher} states the Rademacher expansion, and Sections~\ref{sec:polarinf} to~\ref{circ:sec} prove it: the polar data at the two cusps, the transformation law and the circle method. Section~\ref{sec:model} proves the eta-quotient picture, Section~\ref{sec:bessel} collects the facts on Bessel functions, Section~\ref{sec:errors} bounds every remaining term, and Section~\ref{sec:proofall} proves Theorem~\ref{thm:all}. Section~\ref{sec:values} compares the values with the product formula and with Sen's table, and Section~\ref{sec:remains} describes the extension to the other CHL models. The appendices describe the exact computation and the interval certificates and list the margins.
